\section{Waiting-Room Mechanisms}
\label{sec:design}

All mechanisms are in ResultBench's waiting room
(\texttt{internal/wroom}), which sits in front of the portal and admits
students at a configured rate in first-come-first-served or lottery order. A
student's ticket is kept either in the page (\emph{naive}) or in a signed
cookie (\emph{token}).

\textbf{Admission lease.} An admission is good for at most $k$ requests to
the portal, as well as for a session lifetime (60\,s). A student who has used
$k$ requests and asks again rejoins the queue, by default at the front, so
that a student who arrived early keeps their place but cannot send more than
$k$ requests per admission. The portal then receives at most $k$ times the
admission rate. We test $k \in \{1, 2, 4\}$ and re-entry at the front or the
back.

\textbf{Closed-loop admission.} Instead of a fixed rate, the room measures
every request it forwards. Each second, if more than 10\% of the requests that
finished took longer than a target (500\,ms), failed, or were abandoned by the
student, or if none finished while some were in progress, the rate is
multiplied by 0.8; otherwise, if students are waiting, it grows by 5\% of
the starting rate. The rate stays between 0.1 and 4 times its starting value.
This is additive increase and multiplicative decrease, as in TCP congestion
control and in Breakwater~\cite{cho2020breakwater}, applied to admissions
rather than requests.

\textbf{Wait estimate.} Each queued response carries the student's position
and an estimate of the remaining wait: position divided by the current
admission rate, plus the time to release if the student arrived early. With
the lottery order the position before release is unknown and no estimate is
given.
